\documentclass[../Main.tex]{subfiles}

\begin{document}
\section{Defining a Model}
We have data $X \in \R^{n \times p}$ and $\vec{Y} \in \R^n$. In a linear model we assume:
\begin{equation*}
    \vec{Y} = X\vec{\beta} + \vec{\epsilon}
\end{equation*}
where $\vec{\beta} \in \R^p$ is a vector of coefficients to be determined, and $\vec{\epsilon} \in \R^n$ is a vector of independent and identically distributed random noise variables.

Initially we assume that $\epsilon_i$ have mean $0$ and variance $\sigma^2$ which is unknown. We also assume that $X$ has rank $p$ (full rank).
\section{Least Squares Estimation}
\begin{definition}{Least squares estimator}
    The \underline{least squares estimator} $\hat{\vec{\beta}}$ is the minimiser of the residual sum of squares,
    \begin{equation*}
        S(\vec{\beta}) = \norm{\vec{Y} - X \vec{\beta}}_2^2
    \end{equation*}
    and is given by $\hat{\vec{\beta}} = (X^TX)^{-1} X^T \vec{Y}$.
\end{definition}
\begin{theorem}[Gauss-Markov theorem]
    For an unbiased estimator $\tilde{\vec{\beta}} = C\vec{Y}$ for some matrix $C$,
    \begin{equation*}
        \Var(\vec{t} \cdot \hat{\vec{\beta}}) \leq \Var(\vec{t} \cdot \hat{\vec{\beta}})
    \end{equation*}
    for any $\vec{t} \in \R^p$. We say that $\hat{\vec{\beta}}$ is BLUE (best linear unbiased estimator).
    \label{thmGaussMarkov}
\end{theorem}
\begin{proof}
    First note that the condition is equivalent to the matrix $\Cov(\tilde{\vec{\beta}}) - \Cov(\hat{\vec{\beta}})$ being positive semi-definite.

    Define $A = C - (X^TX)^{-1} X^T$. Note that $\E{A\vec{Y}} = \E[\tilde{\vec{\beta}}] - \E[\hat{\vec{\beta}}] = 0$, and show that this implies $AX = 0$.

    Now expand $\Cov(\tilde{\vec{\beta}})$ to get $\sigma^2 AA^T + \Cov(\hat{\vec{\beta}})$.
\end{proof}
\begin{definition}{Orthogonal projection}
    A matrix $P \in \R^{n \times n}$ is an \underline{orthogonal projection} if it is idempotent $P^2 = P$ and symmetric $P^T = P$.
\end{definition}
\begin{itemize}
    \item This is equivalent to all $\vec{v}$ in the column space of $P$ being unaffected by $P$ ($P\vec{v} = \vec{v}$), and all $\vec{w}$ in the orthogonal complement to the column space being sent to $\vec0$.
    \item If $P$ is an orthogonal projection, so is $I - P$.
    \item We can split norms: $\norm{\vec{Y}}^2 = \norm{P\vec{Y}}^2 + \norm{(I - P)\vec{Y}}^2$.
    \item The eigenvalues of $P$ are all zero or 1, and the trace of $P$ is the rank of $P$.
\end{itemize}
\begin{definition}{Fitted values}
    The \underline{fitted values} in a linear regression are $\hat{\vec{Y}} = X\hat{\vec{\beta}}$ (predicted mean responses). Then the \underline{residuals} are $\vec{Y} - \hat{\vec{Y}}$.
\end{definition}
Letting $P = X(X^TX)^{-1}X^T$, so $\hat{\vec{Y}} = P\vec{Y}$, we see that this $P$ is an orthogonal projection. The column spaces of $P$ and $X$ are the same.
\section{Maximum Likelihood Estimation}
Require now that $\vec{\epsilon} \sim N(\vec0, \sigma^2 I)$. We can derive the MLE and show that it is the same as the least squares estimator $\hat{\vec{\beta}}$. The estimator for the noise variance is:
\begin{equation*}
    \hat{\sigma}^2 = \frac{\norm{\vec{Y} - X\vec{\beta}}^2}{n} = \frac{\norm{(I-P)\vec{Y}}^2}{n}
\end{equation*}
By considering orthogonal projections,
\begin{itemize}
    \item $\hat{\vec{\beta}} \sim N(\vec{\beta}, \sigma^2 (X^TX)^{-1})$,
    \item $\frac{\hat{\sigma^2}}{\sigma^2} n \sim \chi_{n-p}^2$,
    \item $\hat{\vec{\beta}}$ and $\hat{\sigma}^2$ are independent.
\end{itemize}
\subsection{Testing a Single Coefficient}
First note that:
\begin{equation*}
    \frac{\hat{\beta}_1 - \beta_1}{\sqrt{\sigma^2(X^TX)^{-1}_{11}}} \sim N(0, 1)
\end{equation*}
and so we can construct the confidence inverval based on a $t$ test. This leads us to a hypothesis test of:
\begin{align*}
    H_0&: \beta_1 = 0 \\
    H_1 &: \beta_1 \neq 0
\end{align*}
and we reject $H_0$ when:
\begin{equation*}
    \abs{\hat{\beta}_1} > t_{n-p}\left(\frac{\alpha}{2}\right) \sqrt{\frac{(X^TX)^{-1}_{11} \hat{\sigma}^2 n}{n-p}}
\end{equation*}
where $t_{n-p}\left(\frac\alpha2\right)$ is a quantile of a $t$ distribution.
\subsection{Testing a Group of Coefficients}
After possibly re-labelling, we want to test:
\begin{align*}
    H_0&: \beta_1 = \beta_2 = \cdots = \beta_{p_0} = 0 \\
    H_1&: \vec{\beta} \in \R^p
\end{align*}
Now write as block matrices:
\begin{equation*}
    X = \begin{pmatrix}X_0 & X_1\end{pmatrix},\qquad \vec{\beta} = \begin{pmatrix}\vec{\beta^0} \\ \vec{\beta^1}\end{pmatrix}
\end{equation*}
so $X_0 \in \R^{n \times p_0}$ and $\vec{\beta^0} \in \R^{p_0}$ so $H_0$ is $\beta^0 = \vec0$.

Define $P = X(X^TX)^{-1} X^T$ (as normal) and $P_1 = X_1(X_1^TX_1)^{-1} X_1^T$. We can derive that the GLRT for this test is monotonic in:
\begin{equation*}
    \frac{\norm{(I - P_1)\vec{Y}}^2}{\norm{(I - P)\vec{Y}}^2}
\end{equation*}
We need to use that $(I-P)(P - P_1) = 0$ and $P - P_1$ is an orthogonal projection (rank $p_0$), so the test statistic is increasing in:
\begin{equation*}
    F = \frac{\norm{(P - P_1)\vec{Y}}^2}{\norm{(I - P)\vec{Y}}^2} \frac{1 / p_0}{1 / (n-p)}
\end{equation*}
which is an $F$ distribution.
\end{document}